\chapter{Introduction}
\label{chap:introduction}


\section{Motivation}
\label{sec:motivation}

Whether Earth-like conditions are common in the galaxy is one of the
central questions of modern astronomy, and one that observational
data alone has not been able to answer.  Two decades of transit and
radial-velocity surveys have shown that small planets are ubiquitous
around all types of dwarf stars \parencite{Fressin2013,Bergsten2022},
with an occurrence-rate distribution that is bimodal across the
well-studied radius valley near $1.8\,\Rearth$ separating rocky
super-Earths from sub-Neptunes \parencite{Fulton2017}, and that
$\eta_\oplus\sim 0.2$--$0.6$ rocky planets sit in the classically-defined
habitable zone of each solar-like host
\parencite{Bryson2021,Kunimoto2020}.
Yet the observational sample of planets receiving near-Earth
irradiation remains sparse, biased toward large radii, and --- with
the exception of a small number of bright nearby M dwarfs ---
dominated by hosts too faint for the atmospheric follow-up that a
habitability claim would require
\parencite{Bean2017,LustigYaeger2019}.  The habitable zone therefore
remains, more than three decades after its modern formulation
\parencite{Kasting1988,Kopparapu2013}, a theoretical construct
grounded on essentially a single data point.

The clearest illustration of why this matters comes from Venus and
Earth.  The two planets accreted from overlapping mass reservoirs
\parencite{Kleine2020,Raymond2020}, exhibit bulk densities within a
few percent of each other \parencite{Smrekar2018}, and are thought to
have followed broadly similar formation sequences, including a hot
magma-ocean stage driven by the late giant impacts of accretion
\parencite{Gillmann2020,ElkinsTanton2012,Lichtenberg2022}.  Their
present states could not be more different.  Whether Venus lost its
water reservoir late in its evolution, through a slow desiccation
after billions of years of temperate climate \parencite{WayGenio2020},
or already during its magma-ocean phase, through the trapping of a
steam-dominated atmosphere over molten silicate
\parencite{Hamano2013,Turbet2021}, remains open.  The two scenarios
predict radically different geophysical histories --- and radically
different observational signatures on exoplanets caught mid-way
between them.

The runaway greenhouse effect is what separates them, and its
physics is a simple positive feedback.  Water vapour is a strong
greenhouse gas, and a warmer atmosphere holds exponentially more of
it, so any wet planet whose insolation is nudged upward warms
further --- and the further it warms, the more opaque its atmosphere
becomes to its own outgoing thermal radiation.  There is a ceiling
on the rate at which such an atmosphere can radiate to space --- the
Simpson--Nakajima limit \parencite{Nakajima1992,Goldblatt2013} ---
and once the absorbed stellar flux exceeds this ceiling, no
equilibrium with a liquid ocean is possible: the ocean evaporates,
the surface climbs past silicate-melting temperatures, and the
planet settles into a steam-atmosphere-over-magma-ocean state whose
transit signature is measurably different from the temperate one it
replaced.  The numerical value of the threshold, its dependence on
stellar spectrum and planetary albedo, and the size of the resulting
transit-radius inflation are treated in
\Cref{chap:exoplanet_characterization}.

This observation --- that the runaway transition is imprinted on the
transit radius --- reframes the habitable-zone question.  Instead of
attempting to prove that any single planet is habitable, one can ask
whether the \emph{population} of small planets shows the step in
mean transit radius that the runaway greenhouse hypothesis predicts.
A statistical detection of this discontinuity would constitute the
first observational test of the habitable-zone concept itself, and a
non-detection would place quantitative upper limits on the water
inventory of rocky worlds
\parencite{Turbet2019,BixelApai2020,Bean2017}.
\textcite{Schlecker2024} formalised this idea into a concrete
predictive framework, which they termed the habitable-zone inner-edge
discontinuity, or \HZIED.  Their principal finding was that a
demographic detection is within reach of the ESA \PLATO mission,
provided a set of favourable assumptions on stellar age, planetary
water content, and target selection all hold.  Whether that
conclusion survives contact with the \emph{actual} \PLATO Input
Catalog, with the real distribution of stellar ages measured by
Gaia, and with a realistic population of planet compositions
inherited from modern formation models is the question this thesis
takes up.

The current transit catalogue --- assembled over two decades by
\Kepler{}, K2, and \TESS{} --- is deeply informative about what
happens on the hot side of the runaway threshold, and essentially
mute about the cold side.  Of $4\,044$ confirmed \Kepler+K2+\TESS{}
planets with both a measured radius and a computed instellation,
only eight sit below $S_\mathrm{thresh}=280\,\mathrm{W\,m^{-2}}$ at
any radius, and only one satisfies the mildest \HZIED-compatible
cut used elsewhere in this thesis.  The two-mean Bayes factor
introduced as the simplest test of the discontinuity is therefore
formally undefined on today's data
(\Cref{chap:kepler_tess}).  \PLATO exists precisely to change that.


\section{Aims and contributions of this thesis}
\label{sec:contributions}

This thesis extends the \HZIED framework of \textcite{Schlecker2024}
along four axes chosen for their expected impact on the predicted
detectability, and evaluates the result on the actual \PLATO target
samples rather than on a generic Gaia-based synthetic catalogue.

\emph{Real target samples.}  The population generator ingests the
\PLATO Input Catalog directly and produces planet catalogues for
the P4 (M-dwarf, $V\leq 16$) and P5 (FGK, $V\leq 13$) samples defined
by \textcite{Nascimbeni2022}.  Per-sample radius precisions are
adopted from the \PLATO \emph{Definition Study Report}
\parencite{ESA2017PLATO} rather than assumed to be uniform across
the survey.

\emph{Real stellar ages.}  Where the original framework drew stellar
ages from a uniform prior on $[0,\,10]\,\mathrm{Gyr}$, this thesis
uses the Gaia DR3 FLAME isochrone catalogue \parencite{GaiaDR3}
directly, or a bootstrap fill from the FLAME distribution when no
FLAME value is available, or a blend of FLAME plus a uniform prior
for the unmeasured tail.  For age-sensitive steam models such as
\textcite{Aguichine2021}, the choice between these variants shifts
the recovered evidence by tens of $\ln\Delta Z$ units, and in one
case flips the sign of the detection.  The origin of this
sensitivity is quantified in \Cref{ssec:aguichine}.

\emph{Physically extended steam prescriptions.}  Beyond the
\textcite{Turbet2020} grey-atmosphere table that anchored the
original framework, five additional prescriptions are compared side
by side: a 2D AGNI radiative--convective grid, an
AGNI+PALEOS-updated version of that grid, the
\textcite{Aguichine2021} irradiated-ocean model, the
\textcite{Pierrehumbert2022} Hycean prescription, and a mixed
rocky/Hycean population motivated by \textcite{Innes2023}.  The
comparison exposes both the mass range over which each prescription
is defensible and the systematic differences between them at the
fiducial injection.  The mixed model in particular shows that a
composite population with two distinct runaway thresholds can
produce a strong single-step Bayes factor whose recovered
threshold is nowhere near either injected value
(\Cref{ssec:hycean_mixed}) --- a cautionary result for future
interpretation of a real detection.

\emph{Multi-seed noise-realisation treatment.}  Every sweep figure
in this thesis reports the median of ten independent noise
realisations and shades the $16$--$84\,\%$ percentile band across
them.  For typical \PLATO configurations the seed-to-seed spread on
the recovered $\ln\Delta Z$ is comparable in amplitude to the
mode-to-mode systematic across the injected water content, and
single-seed numbers reported in earlier literature therefore need
to be interpreted with the seed-to-seed noise band in mind.  The
same treatment is applied uniformly to the recovered
$\hat{S}_\mathrm{thresh}$, which is used throughout as an
orthogonal diagnostic against the possibility of a decisive
$\ln\Delta Z$ arising from structure the single-step hypothesis
cannot represent.

The overarching question addressed by these extensions is:

\begin{quote}
\emph{Can, and if so, when during mission time can \PLATO detect the population-level radius discontinuity predicted by the runaway greenhouse transition, and which combination of stellar age distribution, planetary composition, and observed target sample is required to make the detection decisive?}
\end{quote}


% \section{Structure of this thesis} removed from condensed version
